% ============================================================================
% all_tables.tex
%
% Assembles every results table of the paper into one standalone document,
% using the table files that are already produced by the pipeline.
%
% Sources (all pulled in with \input so the numbers reflect the latest run):
%   * MATLAB ESS comparison tables (3,4,5,9,10,11,15):
%       diagrams/Table3.txt, Table4.txt, Table5_N64.txt, Table5_N32.txt,
%       diagrams/Table9.txt, Table10.txt, Table11.txt, Table15.txt
%     (columns: Method | Time(s) | Acc. (%) | step size delta | ESS | Min ESS/s)
%   * MATLAB variance-reduction tables (7,8,12):
%       diagrams/Table7.txt, Table8.txt, Table12.txt
%   * R-based tables (2,6,13,14):
%       table2_logistic.tex, Table6.tex, Table13.tex, Table14.tex   (repo root)
%
% Compile AFTER running the experiments and the R scripts:
%   run_all_experiments        % MATLAB, from repo root
%   source("make_table2_logistic.R"); source("make_table6_logistic.R")
%   source("tmcmc.R");          source("tmcmc_ess.R")          % R, from repo root
%   pdflatex all_tables.tex
% ============================================================================
\documentclass[11pt]{article}

\usepackage[margin=1.8cm]{geometry}
\usepackage{booktabs}   % \toprule \midrule \bottomrule
\usepackage{array}      % custom column types
\usepackage{amsmath,amssymb}
\usepackage{float}      % [H] float specifier used by Table14.tex
\usepackage{caption}

% Hidden column: consumes the cell content but prints nothing (zero width).
% Used to drop the "Acc. (%)" and "step size delta" columns (3rd and 4th) from
% the 6-column ESS tables in diagrams/Table*.txt without regenerating them.
\newcolumntype{H}{>{\setbox0=\hbox\bgroup}c<{\egroup}@{}}

\begin{document}

\begin{center}
  {\Large\bfseries Updated result tables}
\end{center}

\bigskip

\noindent\fbox{\parbox{\textwidth}{%
\textbf{Note on the reported numbers.}
The values in the tables below differ from those in the corresponding tables of
\texttt{main\_paper.pdf}. This is because the present results were generated with
a \emph{different random seed}: in the first submission of the paper the random
seed was accidentally not stored, so the exact numerical values reported there
cannot be reproduced. The seed is now fixed at the top of
\texttt{run\_all\_experiments.m}, so all numbers in this document are fully
reproducible. The qualitative conclusions of the paper are unchanged.
}}

\bigskip

% ============================================================================
% Table 2 — ESS, Bayesian logistic regression (R)
% ============================================================================
\begin{table}[htbp]
  \centering
  \caption{ESS scores in Bayesian logistic regression with $d$-dimensional
           posterior distributions. All numbers are averages across ten random
           repeats. (Table 2, Section 5.1.1.)}
  \input{table2_logistic.tex}
\end{table}

% ============================================================================
% Table 3 — Heart (binary classification, MATLAB)
% ============================================================================
\begin{table}[htbp]
  \centering
  \caption{Comparison of sampling methods in the Heart dataset ($d=270$).
           Averages across ten random repeats. (Table 3, Section 5.1.2.)}
  \begin{tabular}{lr@{}HHrr}
    \toprule
    \input{diagrams/Table3.txt}\crcr
    \bottomrule
  \end{tabular}
\end{table}

% ============================================================================
% Table 4 — Australian (binary classification, MATLAB)
% ============================================================================
\begin{table}[htbp]
  \centering
  \caption{Comparison of sampling methods in the Australian Credit dataset
           ($d=690$). Averages across ten random repeats. (Table 4, Section 5.1.2.)}
  \begin{tabular}{lr@{}HHrr}
    \toprule
    \input{diagrams/Table4.txt}\crcr
    \bottomrule
  \end{tabular}
\end{table}

% ============================================================================
% Table 5 — Log-Gaussian Cox process (MATLAB)
% ============================================================================
\begin{table}[htbp]
  \centering
  \caption{Comparison of sampling methods in the log-Gaussian Cox model
           ($d=4096$), $N=64$. Averages across ten random repeats. (Table 5, Section 5.1.3.)}
  \begin{tabular}{lr@{}HHrr}
    \toprule
    \input{diagrams/Table5_N64.txt}\crcr
    \bottomrule
  \end{tabular}
\end{table}

\begin{table}[htbp]
  \centering
  \caption{Comparison of sampling methods in the log-Gaussian Cox model, $N=32$
           (supplementary). Averages across ten random repeats.}
  \begin{tabular}{lr@{}HHrr}
    \toprule
    \input{diagrams/Table5_N32.txt}\crcr
    \bottomrule
  \end{tabular}
\end{table}

% ============================================================================
% Table 6 — VR factors, Bayesian logistic regression (R)
% ============================================================================
\begin{table}[htbp]
  \centering
  \caption{Range of estimated factors (variance ratios) by which the variance of
           $\mu_n(F)$ is larger than the variance of $\mu_{n,G}(F)$ for the
           Bayesian logistic regression targets. (Table 6, Section 5.2.)}
  \input{Table6.tex}
\end{table}

% ============================================================================
% Table 7 — VR factors, GP binary classification, Heart/Australian (MATLAB)
% ============================================================================
\begin{table}[htbp]
  \centering
  \caption{Estimated factors for the Gaussian process binary classification
           targets. (Table 7, Section 5.2.)}
  \begin{tabular}{lrrrr}
    \toprule
    \input{diagrams/Table7.txt}\crcr
    \bottomrule
  \end{tabular}
\end{table}

% ============================================================================
% Table 8 — VR factors, GP log-Cox model (MATLAB)
% ============================================================================
\begin{table}[htbp]
  \centering
  \caption{Estimated factors for the Gaussian process log-Cox model with
           $d=4096$. (Table 8, Section 5.2.)}
  \begin{tabular}{lrr}
    \toprule
    Dataset & $n = 1{,}000$ & $n = 10{,}000$ \\
    \midrule
    \input{diagrams/Table8.txt}\crcr
    \bottomrule
  \end{tabular}
\end{table}

% ============================================================================
% Table 9 — German (binary classification, MATLAB)
% ============================================================================
\begin{table}[htbp]
  \centering
  \caption{Comparison of sampling methods in the German dataset ($d=1000$).
           Averages across ten random repeats. (Table 9, Appendix A.1.)}
  \begin{tabular}{lr@{}HHrr}
    \toprule
    \input{diagrams/Table9.txt}\crcr
    \bottomrule
  \end{tabular}
\end{table}

% ============================================================================
% Table 10 — Pima (binary classification, MATLAB)
% ============================================================================
\begin{table}[htbp]
  \centering
  \caption{Comparison of sampling methods in the Pima dataset ($d=532$).
           Averages across ten random repeats. (Table 10, Appendix A.1.)}
  \begin{tabular}{lr@{}HHrr}
    \toprule
    \input{diagrams/Table10.txt}\crcr
    \bottomrule
  \end{tabular}
\end{table}

% ============================================================================
% Table 11 — Ripley (binary classification, MATLAB)
% ============================================================================
\begin{table}[htbp]
  \centering
  \caption{Comparison of sampling methods in the Ripley dataset ($d=250$).
           Averages across ten random repeats. (Table 11, Appendix A.1.)}
  \begin{tabular}{lr@{}HHrr}
    \toprule
    \input{diagrams/Table11.txt}\crcr
    \bottomrule
  \end{tabular}
\end{table}

% ============================================================================
% Table 12 — VR factors, GP binary classification, German/Pima/Ripley (MATLAB)
% ============================================================================
\begin{table}[htbp]
  \centering
  \caption{Range of estimated variance reduction factors for Gaussian process
           binary classification targets. (Table 12, Appendix A.1.)}
  \begin{tabular}{lrrrr}
    \toprule
    \input{diagrams/Table12.txt}\crcr
    \bottomrule
  \end{tabular}
\end{table}

% ============================================================================
% Table 13 — VR factors, Student-t tail probability (R)  [two sub-tables N=2, N=5]
% ============================================================================
\begin{table}[htbp]
  \centering
  \caption{Estimated factors (variance ratios) by which the variance of
           $\mu_n(F)$ is larger than the variance of $\mu_{n,G}(F)$ when the
           target is a univariate Student's $t$-distribution with $\nu$ degrees
           of freedom and $F(x)=\mathbb{I}(x>b)$, for different $\nu$, $b$ and
           truncation level $N$ ($N=2$ top, $N=5$ bottom). (Table 13, Appendix A.2.)}
  \input{Table13.tex}
\end{table}

% ============================================================================
% Table 14 — ESS / acceptance / step size, Student-t (R)
% This file already provides its own \begin{table}...\end{table} float.
% ============================================================================
\input{Table14.tex}

% ============================================================================
% Table 15 — GP regression with informative likelihood (MATLAB)
% ============================================================================
\begin{table}[htbp]
  \centering
  \caption{Comparison of sampling methods in the regression dataset with
           $\sigma^2 = 0.01$ and $d=1000$. Averages across ten repeats.
           (Table 15, Appendix A.4.)}
  \begin{tabular}{lr@{}HHrr}
    \toprule
    \input{diagrams/Table15.txt}\crcr
    \bottomrule
  \end{tabular}
\end{table}

\end{document}
